\chapter{Plano de análise}
\label{ap:plano}

Este apêndice detalha cinco pontos que o corpo do projeto apenas enuncia. Nada
aqui é resultado: são regras e estruturas a fixar antes do contato com os dados
posteriores a 2019.

\section{Registros da aplicação aérea: cobertura e classificação}

\textbf{Função de cada fonte.} Os relatórios mensais de operadores ao
Ministério da Agricultura informam a aplicação declarada pelo operador. O
receituário agronômico informa a modalidade prescrita e não comprova, sozinho, a
aplicação realizada. Autos e licenças do órgão ambiental estadual informam
fiscalização. O Censo Agropecuário de 2006 serve à caracterização histórica e
não classifica município e mês de 2010 a 2018.

\textbf{Etapa de viabilidade.} Antes de qualquer regra de decisão, verifica-se o
que cada órgão entregou: período coberto, operadores e municípios presentes,
meses sem documento, e se a ausência de documento é declarada pelo órgão como
ausência de atividade ou como falta de arquivo.

\textbf{Classificação.} Cada município e mês recebe uma de três categorias:
\emph{voo}, quando há aplicação aérea declarada por operador ou constatada em
fiscalização (prescrição aérea sem confirmação é registrada à parte, como
indício);
\emph{sem voo}, apenas quando a ausência é sustentada por cobertura suficiente,
isto é, a fonte cobre o período e os operadores que atuavam na região e não
registra aplicação; e \emph{sem registro}, em todos os demais casos. Lacuna não
é zero. Aplicação intermitente é descrita como tal, pela proporção de meses da
estação de aplicação com voo em cada ano, e não reduzida a uma data de
cessação. Divergência entre fontes é reportada, e não resolvida em favor de uma
delas.

\textbf{Regra.} ``Cessação anterior sem retomada'' exige a categoria \emph{sem
voo} em todas as estações de aplicação de um ano em diante até 2018. Os
critérios de cobertura suficiente são fixados depois da etapa de viabilidade e
antes da leitura do conteúdo dos registros por município.

\textbf{Estudo complementar.} Se a cessação anterior for documentada, um estudo
da interrupção teria objeto distinto, a interrupção e não a lei, e
identificação própria: partilhar o clima não garante que a seca tenha afetado
de modo igual a produção, o emprego e a renda dos dois lados, que são caminhos
concorrentes da Figura~\ref{fig:dag}. Esse estudo não substitui a avaliação da
Lei nº~16.820/2019.

\section{Grupo de comparação}

A delimitação proposta para a Região Integrada de Desenvolvimento é ponto de
partida, por ser uma lista externa ao pesquisador. A lista final de controles é
justificada por critérios territoriais e produtivos, fixados antes dos
desfechos: pertencer à mesma chapada ou ao seu entorno imediato; e
características de linha de base (população, urbanização, estrutura agrária,
presença de fruticultura irrigada, rede de saúde). Reporta-se a tabela de
equilíbrio de linha de base entre tratados e controles e uma versão do
contraste restrita aos municípios potiguares com fruticultura irrigada, presumivelmente mais expostos aos mesmos choques econômicos e fitossanitários do setor. Quando os
registros permitirem, as localidades são classificadas pelo grau de exposição
documentado, para descrição e heterogeneidade, sem alterar a definição de
tratados e controles.

\section{Inferência: o que será verificado}

A unidade é o município. Mais nascimentos tornam mais precisa a média
municipal, e não acrescentam realizações independentes de choques municipais ou
estaduais. Se o erro tem um componente comum aos municípios de cada estado,
$u_{it} = a_{s(i),t} + e_{it}$, o contraste entre Ceará e Rio Grande do Norte
carrega a diferença entre os componentes estaduais; placebos formados só entre
municípios potiguares cancelam o componente potiguar e não reproduzem essa
diferença.

Três objetos são tratados em separado: os $ATT(t)$ trimestrais, o efeito
agregado do pós-período e, se usadas, as bandas simultâneas. Cada procedimento
candidato é examinado para o mesmo parâmetro das Seções~\ref{sec:parametro}
e~\ref{sec:estimador}, com a
mesma ponderação por nascimentos e a mesma estrutura temporal, quanto a: como os
controles representam a distribuição dos choques dos tratados; que diferenças
de variância admite; como trata dependência serial e entre municípios; e se é
compatível com a versão duplamente robusta e com covariáveis.

Os exercícios no pré-período atribuem intervenções fictícias preservando as
trajetórias municipais, o porte dos municípios e uma dependência territorial
plausível, sem permutar observações isoladas. São diagnóstico de calibração, e
não validação das hipóteses: bom desempenho entre controles potiguares não
elimina a possibilidade de um choque diferencial entre os estados. Se essa
fonte de incerteza não puder ser representada pelos dados de controle, os
intervalos são apresentados como condicionais a uma restrição explícita sobre
ela, ou acompanhados de sensibilidade a cenários de choque estadual. A mesma
estrutura serve para simular o poder dos parâmetros efetivamente estimados.

\section{Cenários de exposição e efeito mínimo detectável}

O efeito municipal hipotético é o produto da fração de nascimentos expostos pelo
efeito por nascimento exposto, que por sua vez depende da intensidade e da
duração da exposição. Mudar essas hipóteses altera o efeito do cenário, e não o
efeito mínimo detectável, que depende apenas dos tamanhos amostrais e dos
parâmetros de variância. A Tabela~\ref{tab:cenarios} compara os efeitos de
cenário com os mesmos limites da Tabela~\ref{tab:emd}.

\begin{table}[ht]
\centering\small
\caption{Efeito municipal hipotético sobre o peso ao nascer, por cenário (gramas)}
\label{tab:cenarios}
\begin{tabular}{@{}lrrr@{}}
\toprule
 & \multicolumn{3}{c}{Efeito por nascimento exposto} \\
\cmidrule(l){2-4}
Fração de nascimentos expostos & 40 g & 80 g & 160 g \\
\midrule
0,8\% & 0,3 & 0,6 & 1,3 \\
1,4\% & 0,6 & 1,1 & 2,2 \\
2,8\% & 1,1 & 2,2 & 4,5 \\
5\%   & 2,0 & 4,0 & 8,0 \\
10\%  & 4,0 & 8,0 & 16,0 \\
25\%  & 10,0 & 20,0 & 40,0 \\
\bottomrule
\end{tabular}
\fonte{Elaboração própria. O valor de 80 gramas é o de
\citeonline{calzada2023bananas}; 40 e 160 gramas são hipóteses, a metade e o
dobro. As frações de 1,4\% e 2,8\% vêm da Seção~\ref{sec:poder}; 0,8\%
corresponde a dezesseis nascimentos expostos em 2.078; as demais são hipóteses. Limites inferiores do efeito mínimo detectável, inalterados: 23,9 e
18,0 gramas (nove anos de base) e 49,0 e 36,9 gramas (um ano de base).}
\end{table}

Com 80 gramas por nascimento exposto, o efeito de cenário só alcança o menor
limite, de 18,0 gramas, se a fração exposta chegar a 22,5\%. A grade não fixa um
teto para o efeito total da lei: o estudo equatoriano não distingue gradientes
além de 150 metros naquele contexto, o que não determina o alcance da deriva
para outras substâncias e condições, e a razão entre nascimentos expostos e área
plantada no Equador não determina a da Chapada. A fração exposta local é objeto
de medida.

\section{Estrutura do custo-benefício parcial}

\textbf{Perspectiva, contrafactual e horizonte.} A análise é feita da
perspectiva social, contra a situação sem a lei, no horizonte do pós-período
observado.

\textbf{Saúde, desfechos diretamente valoráveis.} Variação em intoxicações e
internações, com o sinal estimado (benefício ou custo), valorada pelo custo da
doença. Cada episódio é contado uma vez: a internação
que decorre de uma intoxicação notificada não se soma a ela.

\textbf{Peso ao nascer.} Um ganho em gramas só tem valor monetário por meio de
uma função em três elos: da variação do peso para a frequência de consequências
clínicas (baixo peso, prematuridade, internação neonatal); dessas consequências
para custos de tratamento e perdas futuras; e desses custos para valores
evitados. Cada elo exige parâmetros externos, com fonte e incerteza declaradas.
Sem essa calibração não há limiar em gramas calculável, e a pergunta invertida
não é apresentada.

\textbf{Custos.} O custo real é a diferença de custo entre as tecnologias de
aplicação. Variações de produção bruta e de emprego não equivalem a perda de
excedente social: são reportadas à parte, como efeitos distributivos e
ocupacionais, e não se somam ao custo de substituição, para não contar duas
vezes o mesmo ajuste.

\textbf{Incerteza.} Os intervalos do primeiro ensaio são propagados aos
resultados, e cada parâmetro de cenário é apresentado com a sua faixa.
